\documentclass[10pt,a4paper]{amsart}
\usepackage[margin=19mm]{geometry}
\usepackage{amsmath,amssymb,mathtools}
\usepackage[scr=rsfs,cal=euler]{mathalfa}
\usepackage{xfrac}
\usepackage[unicode,hidelinks,bookmarks=false]{hyperref}
\newcommand{\ie}{i.e.\ }
\newcommand{\cf}{cf.\ }
\newcommand{\resp}{respectively\ }
\newcommand{\Subsection}{\subsection}
\providecommand{\nobreakdash}{\nobreak\hbox{-}}
\setlength{\emergencystretch}{2em}
\multlinegap0pt
\usepackage{bm}
\usepackage{bbm}
\usepackage{dsfont}
\usepackage{xspace}
\usepackage{cancel}
\usepackage{mathtools}
\usepackage{cleveref}
\usepackage[shortlabels]{enumitem}
\usepackage{amsthm}
\makeatletter
\@ifundefined{theorem}{\newtheorem{theorem}{Theorem}}{}
\@ifundefined{proposition}{\newtheorem{proposition}[theorem]{Proposition}}{}
\makeatother
\newcommand{\betageo}{\beta^{\mathrm{geo}}}
\newcommand{\vpp}[2]{v_{#1}(#2)}
\title[Prime-Exponent Transition Geometry and Divisor Barriers Betw]{Prime-Exponent Transition Geometry and Divisor Barriers Between Consecutive Highly Composite Numbers}
\author{Marco Mantovanelli}
\date{Source: arXiv:2608.17045v1 (CC BY 4.0)}
\begin{document}
\maketitle

\noindent\emph{Source and licence.} Prime-Exponent Transition Geometry and Divisor Barriers Between Consecutive Highly Composite Numbers. Version arXiv:2608.17045v1 (CC BY 4.0), distributed under the Creative Commons Attribution 4.0 International licence, as recorded in the arXiv licence metadata for this exact version and shown on the abstract page https://arxiv.org/abs/2608.17045v1. Full rights notice: licenses/arxiv-2608.17045-CC-BY-4.0.txt.\par\smallskip

\noindent\textit{Editorial scope.} Counted unit: the Theorem whose source label is \texttt{thm:compensated-matching} in the article (source0.tex lines 1024--1065, source label \texttt{thm:compensated-matching}), delivered together with the article's own complete original proof of it (source0.tex lines 1067--1090, one \texttt{proof} environment of 24 lines). No mathematics is written, repaired or re-derived: statement and proof are the article's own text line for line. Required context retained uncounted: prop:mandatory-deletion-upper (lines 610--628). Every definition, display and label of those lines is retained unchanged. Omitted from this package and not redistributed: 1--609, 629--1023, 1066--1066, 1091--1747. Nothing inside a retained excerpt is omitted or altered: every retained body is delivered exactly as the article's lines are written, so no omission receipt of any kind accompanies this package. Citations: the retained excerpts carry no citation command at all, so no bibliography entry of the article is transcribed and no reference is redistributed. Reference adaptation: none was needed and none was made; every reference inside the delivered unit resolves inside it, so no unresolved reference or citation is produced. Printed slips kept literal: no printed word, symbol, hypothesis or number of the counted statement or of its proof was altered. Layout and class: the article's own class is \texttt{article} with packages including fontenc, inputenc, amsmath, amsthm, mathtools, bm, newtxtext, newtxmath, microtype, geometry, booktabs, tabularx, array, enumitem; the wrapper is the lane's pinned \texttt{amsart} editorial wrapper at 10pt on A4, which restates the theorem-like environments the retained text needs and adds the article's own macro definitions that the retained text needs (\texttt{\textbackslash betageo}, \texttt{\textbackslash vpp}), with extra wrapper packages (bm, bbm, dsfont, xspace, cancel, mathtools, cleveref, enumitem). The delivered numbering is the unit's own. Assets: no retained span contains a figure environment, a graphics-inclusion command, a PDF-page inclusion, a table or a TikZ picture; no image file is copied into this package, the package declares no asset dependency and no image file is redistributed with it. Licence: version arXiv:2608.17045v1 of this article is distributed under the Creative Commons Attribution 4.0 International licence, as recorded in the arXiv licence metadata for this exact version and shown on the abstract page https://arxiv.org/abs/2608.17045v1. A full rights notice is delivered at licenses/arxiv-2608.17045-CC-BY-4.0.txt. Characters: every retained excerpt is pure ASCII, so the delivered engine, XeLaTeX with its default Latin Modern text font, reports no missing character.
\begin{proposition}[Mandatory-deletion upper bound]
\label{prop:mandatory-deletion-upper}
Let
\[
 H=\prod p^{a_p}<H'=\prod p^{b_p}
\]
be consecutive highly composite numbers, and suppose that at least one coordinate decreases.  Then
\begin{equation}
 \boxed{
 \betageo(H,H')\le
 \min_{p:a_p>b_p}\frac{b_p+1}{b_p+2}.}
 \label{eq:mandatory-deletion-upper}
\end{equation}
In particular, if a prime disappears from the support, so that $a_p>0=b_p$, then
\begin{equation}
 \betageo(H,H')\le\frac12.
 \label{eq:support-loss-upper}
\end{equation}
\end{proposition}

\begin{theorem}[Compensated deletion matching]
\label{thm:compensated-matching}
Let
\[
 H=\prod_p p^{a_p}<H'=\prod_p p^{b_p}
\]
be consecutive highly composite numbers.  Suppose exactly one source support prime disappears: for some prime $r$,
\[
 a_r=1,
 \qquad
 b_r=0,
\]
and no prime $p\ne r$ satisfies $a_p>0=b_p$.  Delete the layer $(r,1)$ first.  List all remaining required deletion layers as
\[
 D_i:(p_i,A_i),\qquad A_i\to A_i-1,
 \qquad 1\le i\le k,
\]
in a stack-admissible order; explicitly, if $p_i=p_h$ and $i<h$, then $A_i>A_h$.  Match them injectively to distinct required insertion layers
\[
 I_i:(q_i,j_i),\qquad j_i-1\to j_i.
\]
Set $x_0=H/r$.  At stage $i$, put $e_i=\vpp{q_i}{x_{i-1}}$.  Require $e_i<j_i$, insert successively the still missing prerequisite layers
\[
 (q_i,e_i+1),\ldots,(q_i,j_i-1)
\]
followed by $I_i=(q_i,j_i)$, and denote the resulting state by
\[
 y_i=x_{i-1}q_i^{j_i-e_i}.
\]
Then execute $D_i$ and put $x_i=y_i/p_i$.  Suppose, for every $i$, that
\begin{enumerate}[label=\textup{(\alph*)}]
\item none of the prerequisite layers is one of the later matched layers $I_{i+1},\ldots,I_k$;
\item $y_i\le H$;
\item $j_i\le A_i$.
\end{enumerate}
After stage $k$, execute all remaining required insertion layers in any stack-admissible order.  Then this schedule is an $H'$-admissible geodesic and
\begin{equation}
 \boxed{\betageo(H,H')=\frac12.}
 \label{eq:matching-equality}
\end{equation}
If $q_i>p_i$, condition (c) follows automatically from the non-increasing target exponent profile.
\end{theorem}

\begin{proof}
The first deletion changes the $r$-exponent from $1$ to $0$, and therefore
\[
 d(x_0)=\frac12d(H).
\]
The deletion order is stack-admissible, and an insertion coordinate cannot be a deletion coordinate; hence $D_i$ is available when it is called.  At stage $i$, all prerequisite insertions increase both the integer and its divisor count.  They end at the state immediately preceding $I_i$, while $y_i$ is the state immediately after $I_i$.  Consequently every state in this insertion block is at most $y_i\le H$.  Also $x_i=y_i/p_i<y_i$, so the state following $D_i$ respects the same ceiling.

Let $u_i$ be the state immediately before $I_i$.  The prerequisite insertions give $d(u_i)\ge d(x_{i-1})$.  Since $I_i$ raises level $j_i-1$ to $j_i$ and $D_i$ lowers level $A_i$ to $A_i-1$,
\[
 d(x_i)
 =d(u_i)\frac{j_i+1}{j_i}\frac{A_i}{A_i+1}
 \ge d(u_i)
 \ge d(x_{i-1}),
\]
where the first inequality is condition~\textup{(c)}.  Induction therefore shows that every state through the last deletion has at least $d(H)/2$ divisors.

Once all deletions have been performed, the current state divides $H'$.  Every subsequent state obtained by inserting a remaining required layer is again a divisor of $H'$, hence is at most $H'$, and its divisor count can only increase.  The completed schedule is thus an $H'$-admissible geodesic with capacity at least $1/2$.  The deletion $r^1\to r^0$ gives the reverse inequality by \Cref{prop:mandatory-deletion-upper}, proving equality.

If $q_i>p_i$, target monotonicity gives
\[
 j_i\le b_{q_i}\le b_{p_i}\le A_i-1<A_i.
\]
Thus condition~\textup{(c)} is automatic.
\end{proof}

\end{document}
